\newcommand{\DoNotLoadEpstopdf}{}
\documentclass[11pt]{article}
\usepackage[T1]{fontenc}
\usepackage{lmodern}
\usepackage[margin=1.08in]{geometry}
\usepackage{amsmath,amssymb,amsthm,mathtools}
\usepackage{microtype}
\usepackage{xcolor}
\usepackage[colorlinks=true,linkcolor=blue!45!black,citecolor=blue!45!black,urlcolor=blue!45!black]{hyperref}
\input{glyphtounicode}
\input{glyphtounicode-cmex}
\pdfgentounicode=1
\pdftrailerid{}
\pdfsuppressptexinfo=15
\hypersetup{pdftitle={The Dimension of the Two-Adic Hecke Algebra at Odd Level},pdfauthor={OpenAI}}
\newtheorem{theorem}{Theorem}[section]
\newtheorem{proposition}[theorem]{Proposition}
\newtheorem{lemma}[theorem]{Lemma}
\newtheorem{corollary}[theorem]{Corollary}
\theoremstyle{remark}
\newtheorem{remark}[theorem]{Remark}
\numberwithin{equation}{section}
\newcommand{\Z}{\mathbb Z}
\newcommand{\Q}{\mathbb Q}
\newcommand{\C}{\mathbb C}
\newcommand{\F}{\mathbb F}
\newcommand{\OO}{\mathcal O}
\newcommand{\frakx}{\mathfrak x}
\newcommand{\frakm}{\mathfrak m}
\newcommand{\eps}{\varepsilon}
\DeclareMathOperator{\Spec}{Spec}
\DeclareMathOperator{\End}{End}
\DeclareMathOperator{\Hom}{Hom}
\DeclareMathOperator{\GL}{GL}
\DeclareMathOperator{\tr}{tr}
\DeclareMathOperator{\ad}{ad}
\DeclareMathOperator{\Frob}{Frob}
\DeclareMathOperator{\Frac}{Frac}
\DeclareMathOperator{\Fil}{Fil}
\DeclareMathOperator{\height}{ht}
\DeclareMathOperator{\edim}{edim}
\DeclareMathOperator{\Der}{Der}
\DeclareMathOperator{\Maps}{Maps}
\title{The Dimension of the Two-Adic Hecke Algebra\\at Odd Level}
\author{OpenAI}
\date{October 5, 2026}

\begin{document}
\maketitle

\begin{abstract}
For every odd positive integer $N$, every irreducible component of the
full two-adic Hecke algebra of level $\Gamma_1(N)$ has Krull dimension
four. This proves the $p=2$ case of Emerton's dimension conjecture,
including all residual components.
\end{abstract}

\section{Introduction}
\label{sec:introduction}

The congruences between classical modular forms of different weights
assemble their Hecke eigenvalues into a single $p$-adic algebra. Its
spectrum is larger than an individual family of eigenforms: many such
families can pass through the same classical point. The dimension of
this spectrum measures how much variation is present when no slope or
residual irreducibility condition is imposed.

Fix an odd positive integer $N$. For $i\geq1$, let $M_i(N)$ be the
complex vector space of modular forms of weight $i$ and level
$\Gamma_1(N)$, including Eisenstein series. For every prime $\ell\nmid N$,
write $T_\ell$ for the usual Hecke operator and let $S_\ell$ act on
$M_i(N)$ as $\ell^{i-2}\langle\ell\rangle$, where
$\langle\ell\rangle$ is the diamond operator. Let
\[
 T_{\leq k}^{(2)}(N)
 \subseteq \End_{\C}\!\left(\bigoplus_{i=1}^k M_i(N)\right)
\]
be the $\Z$-algebra generated by $T_\ell$ and $\ell S_\ell$ for
$\ell\nmid2N$. Set
\begin{equation}\label{eq:algebra}
 A_k=\Z_2\otimes_{\Z}T_{\leq k}^{(2)}(N),
 \qquad A=T_2(N)=\varprojlim_k A_k,
\end{equation}
with transition maps given by restriction. The topology on $A$ is the
inverse-limit topology, with the $2$-adic topology on each $A_k$.
All dimensions below are Krull dimensions.

\begin{theorem}\label{thm:main}
For every odd positive integer $N$, every irreducible component of
$\Spec T_2(N)$ has dimension exactly $4$.
\end{theorem}

This is the $p=2$ case of Emerton's Conjecture~2.9
\cite{Emerton11}. The algebra in \eqref{eq:algebra} retains all weights,
all Eisenstein systems, and all residual systems of eigenvalues. In
particular, the theorem does not require a residual Galois representation
to be irreducible, nonscalar, or distinguished at $2$.

\subsection*{History and the main input}
Hida's ordinary families \cite{Hida86} and the eigencurve of Coleman and
Mazur \cite{ColemanMazur98} give fundamental constructions of $p$-adic
families of modular eigenforms. Gouv\^ea and Mazur's infinite fern
\cite{GouveaMazur98} explains how families intersect and produce
higher-dimensional Zariski closures. Emerton's account
\cite{Emerton11} formulates the dimension question for the full
prime-to-$Np$ Hecke algebra, without restricting to one residual
deformation problem. His Corollary~2.28 gives the lower bound needed
here: every irreducible component has dimension at least $4$.

The complementary upper bound is a deformation-theoretic question.
Mazur's deformation theory \cite{Mazur89} relates infinitesimal
deformations of a Galois representation to its adjoint cohomology.
For automorphic representations, the finite local conditions of
Bloch and Kato \cite{BlochKato90} isolate a Selmer group within this
cohomology. Under suitable hypotheses, Kisin \cite{Kisin04} proved
geometric adjoint Selmer vanishing and applied it to local deformation
and eigencurve geometry. Allen \cite{Allen16} established adjoint Selmer vanishing
under residual hypotheses and related local genericity to unobstructed
deformations. Newton and Thorne \cite{NewtonThorne23} subsequently
proved the vanishing theorem used here without the residual restrictions
that would obstruct its application to all components at $2$.

We use their Theorem~5.4 for classical cuspidal forms over $\Q$.
The theorem includes $p=2$ and also includes a CM form provided its
CM field is not contained in the cyclotomic $2$-power extension.
Odd level guarantees this last condition. The result concerns the
\emph{full} adjoint representation, of dimension four, which is the
one required when the determinant is allowed to vary.

\subsection*{Method and organization}
The proof works at characteristic-zero points, even on components
whose residual systems are reducible or scalar. There are three steps.
First, the integral algebra is a finite product of complete Noetherian
local rings, and its classical points are Zariski dense. At $2$ we
justify Noetherianity using Chenevier's determinant laws
\cite{Chenevier14}, avoiding division by $2$ in the integral
pseudorepresentation. Second, the Eisenstein locus has dimension at
most $2$. Together with Emerton's lower bound, this supplies a classical
cuspidal point of weight at least $3$ on every component. Third, the
tangent space at such a point has dimension at most $3$: the
Newton--Thorne theorem annihilates the global finite Selmer group,
and the only nonzero local quotient has dimension $3$ at $2$.
A dimension formula then adds the one-dimensional arithmetic quotient.

Two passages in this argument are useful beyond the immediate
application. Section~\ref{sec:eisenstein} constructs the Eisenstein
interpolation maps directly from the inverse-limit Hecke algebra by a
compact graph argument. Section~\ref{sec:tangent} gives a
characteristic-zero tangent bound for a compact Noetherian algebra
carrying a two-dimensional determinant. In particular, it proves the
continuity of the derivations that occur in its algebraic tangent
space. Neither passage assumes residual irreducibility.

We supply these arguments in detail. The deep external inputs are
the classical attachment and local compatibility of Galois
representations, Emerton's lower bound, Chenevier's determinant
theorems, and Newton--Thorne's Selmer vanishing. Their relevant
forms and hypotheses are stated where they are used.

\section{The integral algebra and its Galois determinant}
\label{sec:integral}

Our first goal is to put the inverse limit \eqref{eq:algebra} in the
category of Noetherian rings while retaining all of its classical
systems. Write $S$ for the set of finite primes dividing $2N$, and
let $G=G_{\Q,S}$ be the Galois group of the maximal extension of $\Q$
unramified at finite primes outside $S$. Ramification at infinity is
allowed. All Galois cohomology will be continuous cohomology.

\subsection{Finite stages and classical systems}
We use the classical rationality, integrality, and simultaneous
semisimplicity of the prime-to-level Hecke action; a convenient
statement including weight one and Eisenstein forms is
\cite[Proposition~1.12 and Remark~1.13]{Emerton11}. For cusp forms,
old copies have the same good-prime eigenvalues as their associated
newforms. On the Eisenstein subspace the eigenvalues are sums and
products of Dirichlet character values and powers of primes.

Consequently, for every nonzero finite stage there is an embedding
\begin{equation}\label{eq:finite-embedding}
 A_k\lhook\joinrel\longrightarrow
 \prod_{j=1}^{m_k}\OO_{E_{k,j}},
\end{equation}
where each $E_{k,j}/\Q_2$ is finite and the coordinate maps are
classical eigenvalue systems, with all necessary $2$-adic embeddings
included. Indeed, before tensoring with $\Z_2$ the algebra embeds
into a finite product of rings of integers of number fields. It is
therefore finite and torsion-free over $\Z$; flatness of $\Z_2$
preserves the embedding. Thus $A_k$ is finite free over $\Z_2$, and
its image in \eqref{eq:finite-embedding} is closed. Zero initial
stages, if present, can be discarded.

A \emph{classical point} of $\Spec A$ means the kernel of a
classical eigenvalue map $A\to\OO_E$ obtained from some finite
stage. The following consequences will be used repeatedly.

\begin{lemma}\label{lem:density}
The ring $A$ is compact, Hausdorff, and reduced. Its classical points
are Zariski dense, and the operators $T_\ell,\ell S_\ell$,
$\ell\nmid2N$, topologically generate it over $\Z_2$.
\end{lemma}
\begin{proof}
The restriction maps between finite-stage algebras are surjective:
the generators at the smaller stage are restrictions of the same
generators at the larger stage. Compactness and Hausdorffness follow
from \eqref{eq:algebra}. Reducedness follows from the injections
\eqref{eq:finite-embedding}. Those injections also show that the
intersection of all classical kernels in $A$ is zero, which is
equivalent to the asserted density. Finally, the $\Z_2$-algebra of
polynomials in the indicated generators surjects onto each $A_k$,
and is therefore dense in the inverse limit.
\end{proof}

For every classical system, the theorems of Deligne and
Deligne--Serre \cite{Deligne71,DeligneSerre74} give a continuous
two-dimensional representation of $G$ over a finite extension of
$\Q_2$; in the Eisenstein case one uses a direct sum of characters.
Choose the Frobenius convention so that its characteristic polynomial
at $\ell\nmid2N$ is
\begin{equation}\label{eq:frobenius}
 Z^2-T_\ell Z+\ell S_\ell
\end{equation}
after specialization. A stable lattice ensures integral traces and
determinants at every element of $G$.

\subsection{An integral determinant law}
A two-dimensional \emph{determinant} over a commutative ring $R$ is
a unital multiplicative homogeneous polynomial law of degree two
$R[G]\to R$, compatible with extension of scalars; see
\cite[Section~1]{Chenevier14}. For a determinant $D$, its trace
and determinant on a group element $g$ are defined by
$D(Z-g)=Z^2-t(g)Z+d(g)$. This notion works over rings
in which $2$ is not invertible.

\begin{lemma}\label{lem:determinant}
There is a continuous two-dimensional determinant over $A$ whose
trace and determinant satisfy
\[
 t(\Frob_\ell)=T_\ell,\qquad
 d(\Frob_\ell)=\ell S_\ell\qquad(\ell\nmid2N).
\]
\end{lemma}
\begin{proof}
At a fixed finite stage, take the tuple of traces and determinants
of its classical Galois representations. The tuple belongs to the
image of \eqref{eq:finite-embedding} on Frobenius conjugacy classes,
by \eqref{eq:frobenius}. The union of these classes is dense in $G$
by Chebotarev. The tuples are continuous class functions, and the
image of \eqref{eq:finite-embedding} is closed. Hence the tuples
belong to $A_k$ at every $g\in G$. Uniqueness on the dense union of
Frobenius classes makes these functions compatible as $k$ varies.
They give continuous functions $t,d$ with values in $A$.

For completeness, the determinant law itself can be constructed
without recovering it by division by $2$. On a finite formal sum
$\sum_g x_g g$, prescribe the quadratic polynomial
\begin{equation}\label{eq:quadratic-law}
 \sum_g d(g)x_g^2+
 \sum_{g<h}\bigl(t(g)t(h)-t(gh)\bigr)x_gx_h,
\end{equation}
where any ordering of the finite support may be used. In every
classical representation this is the determinant of
$\sum_g x_g r(g)$; the formula is symmetric since $t(gh)=t(hg)$.
All polynomial identities expressing normalization and
multiplicativity hold coefficientwise in every classical
specialization. Their coefficients therefore vanish in $A$ by
Lemma~\ref{lem:density}. Formula~\eqref{eq:quadratic-law} defines
the required polynomial law over arbitrary $A$-algebras, with trace
$t$ and determinant $d$.
\end{proof}

\subsection{Noetherianity at two}
\begin{proposition}\label{prop:noetherian}
The ring $A$ is a finite product of complete Noetherian local rings
with finite residue fields. Its given topology on each factor is
the maximal-ideal-adic topology.
\end{proposition}
\begin{proof}
Jochnowitz's finiteness theorem \cite{Jochnowitz82}, in the form
\cite[Proposition~2.8]{Emerton11}, says that the reductions of the
classical systems at fixed prime-to-$2$ level form a finite set.
Every maximal ideal of $A_k$ is detected by such a reduction:
the product in \eqref{eq:finite-embedding} is finite integral over
$A_k$, so one can apply lying over. Each $A_k$ is a finite product
of complete local rings with finite residue fields. A surjective
transition map $A_{k+1}\to A_k$ injects the set of maximal ideals
of $A_k$ into that of $A_{k+1}$. These increasing sets have bounded
cardinality, so stabilize on a tail. The corresponding local-factor
maps are surjective, and identify their residue fields. We obtain
\[
 A=\prod_{\alpha=1}^s A^{(\alpha)},
\]
where each $A^{(\alpha)}$ is a profinite local ring with finite
residue field $k_\alpha$.

The canonical Witt coefficient maps $W(k_\alpha)$ into the
finite-stage local rings are the unique unramified lifts of the
residue-field identifications. Their uniqueness makes them
compatible, and gives $A^{(\alpha)}$ its $W(k_\alpha)$-algebra
structure. Project the determinant of
Lemma~\ref{lem:determinant} to this factor, and let
$\overline D_\alpha$ be its reduction. Chenevier's determinant
deformation theorem \cite[Propositions~3.3 and~3.7]{Chenevier14}
provides a complete Noetherian local universal ring
$R_{\overline D_\alpha}$ with finite residue field, representing
continuous deformations of $\overline D_\alpha$. Here the relevant
finiteness hypothesis is Mazur's condition: every open subgroup
$H\subseteq G$ has only finitely many continuous homomorphisms to
$\Z/2\Z$. For $G_{\Q,S}$ this follows from number-field finiteness
for quadratic extensions with ramification restricted to a fixed
finite set; see also \cite[Example~3.6]{Chenevier14}.

Apply the universal property to finite quotients and pass to the
inverse limit. This gives a continuous map
\[
 R_{\overline D_\alpha}\longrightarrow A^{(\alpha)}.
\]
Its image contains the projections of every $T_\ell$ and
$\ell S_\ell$. The image is compact and therefore closed, so
Lemma~\ref{lem:density} makes the map surjective. Its kernel is
closed; the quotient is a complete Noetherian local ring. The
continuous bijection from this compact quotient to the Hausdorff
ring $A^{(\alpha)}$ is a homeomorphism. This proves the topology
assertion as well.
\end{proof}

The determinant argument supplies the $p=2$ justification for the
pseudodeformation step in \cite[Theorem~2.7]{Emerton11}, whose
original references included an odd-prime restriction. No division
by $2$ has been used in Proposition~\ref{prop:noetherian}.

We now record the geometric input from the infinite fern.
\begin{theorem}[Emerton]\label{thm:lower}
For the algebra $A$ in \eqref{eq:algebra}, every irreducible component
of $\Spec A$ has dimension at least $4$.
\end{theorem}
This is \cite[Corollary~2.28]{Emerton11}, with the definitions in
Sections~2.1 and~2.5 of that paper and with $p=2$. It concerns the
Hecke algebra of all modular forms used here, rather than a
universal deformation ring for an irreducible residual
representation.

\section{Finding cuspidal points on every component}
\label{sec:eisenstein}

The lower bound in Theorem~\ref{thm:lower} lets us exclude small
classical loci. We show that Eisenstein systems cannot fill a
component, even when all weights are allowed.

\begin{proposition}\label{prop:exceptional}
The closure of the classical points of bounded weight has dimension
at most $1$. The closure of all Eisenstein points has dimension at
most $2$.
\end{proposition}
\begin{proof}
Systems of weights at most $k$ factor through $A_k$, which is
finite over $\Z_2$. Their closure therefore lies in the
one-dimensional closed subset defined by $\ker(A\to A_k)$.
In particular, we may discard weights one and two when considering
the Eisenstein locus.

In weight $i\geq3$, an Eisenstein system has eigenvalues
\begin{equation}\label{eq:eisenstein-values}
 T_\ell\longmapsto\psi(\ell)+\phi(\ell)\ell^{i-1},
 \qquad
 \ell S_\ell\longmapsto\psi(\ell)\phi(\ell)\ell^{i-1},
\end{equation}
where $\psi,\phi$ are Dirichlet characters of conductors dividing
$N$. The standard Eisenstein description gives these formulas;
oldforms do not change the good-prime system. There are finitely
many possible pairs, including their $2$-adic embeddings.
Partition the systems by pair and by the parity of $i$. A family
with finitely many weights has already been dealt with, so fix an
infinite family, indexed by a set $I$ of weights, and a finite
extension $E/\Q_2$ containing the character values. Put $\OO=\OO_E$.

Every odd prime has a unique expression
\[
 \ell=s_\ell5^{b_\ell},\qquad
 s_\ell\in\{1,-1\},\quad b_\ell\in\Z_2.
\]
On the fixed parity class, $s_\ell^{i-1}$ is constant. In
\eqref{eq:eisenstein-values}, replace $\ell^{i-1}$ by
$s_\ell^{i-1}(1+X)^{b_\ell}\in\OO[[X]]$. The resulting
series specialize correctly at
\begin{equation}\label{eq:weight-specialization}
 X=x_i:=5^{i-1}-1\in4\Z_2,\qquad i\in I.
\end{equation}
We next verify that this prescription extends from generators to
the inverse-limit algebra.

Give $\OO[[X]]$ its maximal-ideal-adic topology. Let $C$ be the
closure in $A\times\OO[[X]]$ of the $\Z_2$-algebra generated by
the paired Hecke generators and their prescribed series. The
projection $C\to A$ is surjective: its image is closed by
compactness and contains a dense subalgebra. For every $(a,F)\in C$
and $i\in I$, continuity gives
\begin{equation}\label{eq:graph-evaluation}
 \lambda_i(a)=F(x_i),
\end{equation}
where $\lambda_i$ is the corresponding classical system.
If $(0,F)\in C$, then $F$ vanishes at all the distinct $x_i$.
The series $F(4Y)$ is a restricted power series, since its
coefficients tend to zero. Strassmann's theorem therefore implies
$F(4Y)=0$, and hence $F=0$. Thus $C\to A$ is also injective.
It is a homeomorphism, and the second projection defines a
continuous map
\begin{equation}\label{eq:eisenstein-map}
 \Phi:A\longrightarrow\OO[[X]].
\end{equation}

Let $B=\Phi(A)$, a compact, hence closed, subring of $\OO[[X]]$.
Choose a prime $\ell\nmid2N$ with
\[
 \ell\equiv1\pmod N,\qquad \ell\equiv5\pmod8.
\]
Such primes exist by the Chinese remainder theorem and Dirichlet's
theorem. For this prime $\psi(\ell)=\phi(\ell)=1$,
$s_\ell=1$, and $b_\ell\in\Z_2^\times$. It follows from the
image of $\ell S_\ell$ that
$u=(1+X)^{b_\ell}\in B$. The series
\[
 u^{b_\ell^{-1}}=
 \sum_{n\geq0}\binom{b_\ell^{-1}}{n}(u-1)^n=1+X
\]
converges in $\OO[[X]]$, and its partial sums belong to $B$.
Closedness gives $X\in B$ and then $\Z_2[[X]]\subseteq B$.
A finite $\Z_2$-basis of $\OO$ consequently generates
$\OO[[X]]$ as a $B$-module. This is a finite integral extension,
so $\dim B=\dim\OO[[X]]=2$. All the points in this family
factor through $B=A/\ker\Phi$ by \eqref{eq:graph-evaluation};
their closure has dimension at most $2$. Taking the finite union
over the families, and adding weights one and two, proves the
claim.
\end{proof}

\begin{corollary}\label{cor:classical-point}
For every minimal prime $P$ of $A$, there is a classical cuspidal
map $\lambda:A\to\OO_E$ of weight at least $3$ such that
$P\subseteq\frakx:=\ker\lambda$ and $\frakx$ belongs to no other
irreducible component of $\Spec A$.
\end{corollary}
\begin{proof}
Let $Z$ be the closure of the Eisenstein points and of the points
of weights one and two. By Proposition~\ref{prop:exceptional},
$\dim Z\leq2$, whereas $\dim V(P)\geq4$ by
Theorem~\ref{thm:lower}. Since $A$ is Noetherian, it has finitely
many minimal primes. Removing $Z$ and every component other than
$V(P)$ leaves an open subset of $\Spec A$ containing the point
$P$. This open subset is nonempty and lies in $V(P)$.
Lemma~\ref{lem:density} supplies a classical point in it, which
has the asserted properties.
\end{proof}

Fix a point furnished by Corollary~\ref{cor:classical-point}. Its
system comes from a cuspidal newform $f$ of weight $k\geq3$ and
level dividing $N$. Enlarging $E$ if necessary, let
\begin{equation}\label{eq:r}
 r:G\longrightarrow\GL_2(E)
\end{equation}
be its absolutely irreducible Galois representation. Absolute
irreducibility here is a characteristic-zero theorem for cuspidal
newforms \cite[Theorem~2.3]{Ribet77}; nothing is asserted about its
reduction. The
specializations $\lambda(t),\lambda(d)$ are the trace and
determinant of $r$.

\section{A tangent bound in characteristic zero}
\label{sec:tangent}

We prove the tangent estimate in a form that separates it from the
particular Hecke algebra. For a representation $r$, the notation
$\ad r$ means $\End_E(E^2)$ with the conjugation action, not its
trace-zero subspace.

\begin{proposition}\label{prop:tangent}
Let $R$ be a compact Hausdorff Noetherian topological
$\Z_p$-algebra, with continuous structure map $\Z_p\to R$.
Suppose a profinite
group $H$ has a continuous two-dimensional determinant over $R$,
with trace $t$ and determinant $d$, and that their values
topologically generate $R$ over $\Z_p$. Let
$\lambda:R\to\OO_E$ be a continuous $\Z_p$-algebra map, where
$E/\Q_p$ is finite, and suppose its specialized determinant is
that of an absolutely irreducible continuous representation
$r:H\to\GL_2(E)$. For $\frakx=\ker\lambda$,
\begin{equation}\label{eq:tangent-bound}
 \dim R_{\frakx}\leq\edim R_{\frakx}
 \leq\dim_E H^1(H,\ad r).
\end{equation}
\end{proposition}

\begin{proof}
The first inequality is the usual dimension bound for a
Noetherian local ring. We prove the second by relating derivations
to first-order Galois deformations.

\smallskip\noindent
\emph{The algebraic tangent space consists of continuous derivations.}
The image $R/\frakx$ is a $\Z_p$-submodule of $\OO_E$, so is
finite over $\Z_p$. Its fraction field
$E_0=\Frac(R/\frakx)$ is the residue field of $R_{\frakx}$ and
is finite over $\Q_p$. Since $p\notin\frakx$, $R_{\frakx}$ is
a $\Q_p$-algebra. Put $\frakm=\frakx R_{\frakx}$. Separability
of $E_0/\Q_p$ gives a coefficient field in
$R_{\frakx}/\frakm^2$, and every $\Q_p$-derivation kills this
field. It follows that
\begin{equation}\label{eq:derivations}
 \Der_{\Q_p}(R_{\frakx},E)
 \simeq\Hom_{E_0}(\frakm/\frakm^2,E),
\end{equation}
so the left side has $E$-dimension $\edim R_{\frakx}$.
The coefficient field can also be obtained by lifting a primitive
element of the finite separable extension across the square-zero
ideal; its minimal polynomial has invertible derivative.

Let $D$ be such a derivation, restricted to $R$. It is
$\Z_p$-linear and vanishes on $\frakx^2$. Noetherianity makes
$\frakx/\frakx^2$ finite over $R/\frakx$, and the exact sequence
\[
 0\longrightarrow\frakx/\frakx^2\longrightarrow
 R/\frakx^2\longrightarrow R/\frakx\longrightarrow0
\]
shows that $R/\frakx^2$ is finite over $\Z_p$. Every finitely
generated ideal of $R$ is closed: it is the image of a continuous
map $R^m\to R$, hence compact. The quotient topology on
$R/\frakx^2$ is its usual finite-module topology, since a
continuous surjection from a finite free $\Z_p$-module to it is
a quotient map between compact Hausdorff spaces. Any
$\Z_p$-linear map from a finite $\Z_p$-module to $E$ is
continuous. In particular, $D:R\to E$ is continuous.

\smallskip\noindent
\emph{First-order determinants lift to continuous representations.}
Write $E[\eps]$ for the dual-number $E$-algebra, with $\eps^2=0$.
The map
$\lambda+\eps D:R\to E[\eps]$ specializes the determinant to
a continuous determinant lifting that of $r$. Passing to the
canonical Cayley--Hamilton quotient of $E[\eps][H]$ (the quotient
imposing the characteristic-polynomial identities of this determinant),
the absolutely
irreducible lifting theorem for determinants
\cite[Theorem~2.22(i) and Corollary~2.23]{Chenevier14} gives a
representation
\[
 r_D:H\longrightarrow\GL_2(E[\eps])
\]
with this determinant, whose reduction can be identified with $r$.
Its applicability is over the Henselian local ring $E[\eps]$:
the residual representation at this point is the split absolutely
irreducible representation $r$ over the characteristic-zero
field $E$.

Here continuity of $r_D$ follows from continuity of its trace.
Choose $g_1,\ldots,g_4\in H$ such that the matrices $r(g_j)$
form an $E$-basis of $M_2(E)$; absolute irreducibility guarantees
such a choice. The lifts $r_D(g_j)$ form a basis over $E[\eps]$.
The trace pairing on this matrix algebra is nondegenerate, so
the coordinates of $r_D(g)$ in that basis are obtained by
inverting a fixed Gram matrix and using the four functions
\[
 g\longmapsto\tr(r_D(g)r_D(g_j))
       =\lambda(t(gg_j))+\eps D(t(gg_j)).
\]
These functions are continuous. Hence $r_D$ is continuous.

\smallskip\noindent
\emph{Traces detect tangent vectors.}
Write $r_D(g)=(1+\eps c(g))r(g)$. The homomorphism identity
is exactly the cocycle identity for a continuous
$c\in Z^1(H,\ad r)$. Comparing traces gives
\[
 D(t(g))=\tr(c(g)r(g)).
\]
Coboundaries have zero image under the linear map
\begin{equation}\label{eq:cohom-traces}
 H^1(H,\ad r)\longrightarrow\Maps(H,E),\qquad
 [c]\longmapsto\bigl(g\mapsto\tr(c(g)r(g))\bigr),
\end{equation}
where $\Maps(H,E)$ is the $E$-vector space of all functions $H\to E$.
Thus the image of the linear map $D\mapsto D\circ t$ lies in
the image of \eqref{eq:cohom-traces}.

This map on derivations is injective. Indeed the degree-two
trace identity gives
\[
 2d(g)=t(g)^2-t(g^2).
\]
If $D(t(g))=0$ for every $g$, applying $D$ and working in the
characteristic-zero field $E$ gives $D(d(g))=0$ for every $g$.
The derivation then vanishes on the algebra of generators, and
on its closure by continuity. Its restriction to $R$ determines
it on $R_{\frakx}$ by the quotient rule. The injection and
\eqref{eq:derivations} now prove \eqref{eq:tangent-bound}.
\end{proof}

Apply Proposition~\ref{prop:tangent} to $R=A$, $p=2$, $H=G$,
and the classical point chosen above. Propositions and lemmas in
Section~\ref{sec:integral} verify its hypotheses. With $V=\ad r$,
we obtain
\begin{equation}\label{eq:hecke-tangent}
 \dim A_{\frakx}\leq\edim A_{\frakx}
 \leq \dim_E H^1(G,V).
\end{equation}
The use of characteristic zero in the last proof does not impose
any condition on the residual system of the component.

\section{The adjoint cohomology bound}
\label{sec:selmer}

We now show that the right side of \eqref{eq:hecke-tangent} is
at most $3$. For a finite place $v$, write
$H^j(\Q_v,V)$ for the cohomology of its absolute Galois group.
The Bloch--Kato finite subspace is
\[
 H^1_f(\Q_v,V)=
 \begin{cases}
 \ker\bigl(H^1(\Q_v,V)\to H^1(I_v,V)\bigr),&v\ne2,\\
 \ker\bigl(H^1(\Q_2,V)\to
 H^1(\Q_2,V\otimes_{\Q_2}B_{\mathrm{cris}})\bigr),&v=2,
 \end{cases}
\]
where $I_v$ is inertia and $B_{\mathrm{cris}}$ is Fontaine's
crystalline period ring. The global group $H^1_f(\Q,V)$ consists
of classes satisfying these conditions at every finite place.
Conditions at infinity impose nothing here: the positive-degree
cohomology of the real Galois group with $E$-coefficients vanishes,
since $2$ is invertible in $E$.

\subsection{Global Selmer vanishing}
The input from \cite[Theorem~5.4]{NewtonThorne23}, specialized
to $\GL_2/\Q$, is the following: if a regular algebraic cuspidal
automorphic representation is non-CM, or is CM by a quadratic
field not contained in $\Q(\zeta_{2^\infty})$, then the finite
Selmer group of its full $2$-adic adjoint representation vanishes.
There is no hypothesis of residual irreducibility or an odd-prime
restriction in this statement.

\begin{lemma}\label{lem:global-selmer}
For the representation \eqref{eq:r},
$H^1_f(\Q,\ad r)=0$.
\end{lemma}
\begin{proof}
The form $f$ has weight at least $3$, so its associated cuspidal
automorphic representation is regular algebraic, in the usual
cohomological normalization. Twists introduced by normalization
do not change its adjoint representation.

Only the CM case requires a check. If $f$ has CM by $K$, its
automorphic representation has the quadratic self-twist
$\chi_{K/\Q}$. Since the level of $f$ divides the odd integer
$N$, its local representation at $2$ is unramified. The local
Langlands parameter of an unramified representation has trivial
inertia; tensoring it with a ramified character cannot leave it
unchanged. Thus $\chi_{K/\Q}$ is unramified at $2$.
The three quadratic subfields of $\Q(\zeta_{2^\infty})$ are
$\Q(\sqrt{-1})$, $\Q(\sqrt2)$, and $\Q(\sqrt{-2})$,
all ramified at $2$. Therefore $K$ is not contained in that
cyclotomic extension. Newton--Thorne's theorem applies in both
cases and gives the result. We may enlarge the finite coefficient
field to apply the theorem; vanishing descends to $E$ under finite
extension of scalars.
\end{proof}

\subsection{The local quotients}
Global Selmer vanishing will make restriction to the sum of the
local quotients injective. To compute these quotients, we first use
local Tate duality to eliminate the $H^2$ terms from the local
Euler characteristic formulas.

\begin{lemma}\label{lem:local-invariants}
For every finite place $v$,
$H^0(\Q_v,V^*(1))=0$, where $V=\ad r$.
\end{lemma}
\begin{proof}
The trace pairing identifies $V^*$ with $V$. The asserted group
is therefore $\Hom_{G_{\Q_v}}(r,r(1))$.

At $v=2$, the representation $r$ is de Rham, hence Hodge--Tate,
with two weights whose difference is $k-1\geq2$. These are the
classical comparison and filtration statements for modular Galois
representations; see \cite[Section~11.3, especially (11.3.3)--(11.3.4)]{Kato04}.
One may pass to an auxiliary fine level when applying that
cohomological description; the representation attached to $f$ is unchanged.
A Tate twist shifts both weights by one, so the weight sets of
$r$ and $r(1)$ are disjoint. Subrepresentations and quotients of
Hodge--Tate representations are Hodge--Tate, with their weights
among those of the original representation. A nonzero image of
a map $r\to r(1)$ would have weights in both sets, which is
impossible.

Let $v\ne2$. Local--global compatibility for cuspidal newforms
\cite[Theorem~A]{Carayol86} identifies the Frobenius-semisimplified
Weil--Deligne parameter of $r|_{G_{\Q_v}}$ with that of the
local automorphic representation, up to the conventional
normalizing twist. The local automorphic representation is
generic. The genericity criterion
\cite[Lemma~1.1.3]{Allen16} says that its Frobenius-semisimple
Weil--Deligne parameter $W$ has
$\Hom_{\mathrm{WD}}(W,W(1))=0$. Here the Tate twist is normalized
consistently with geometric Frobenius; normalizing character
twists do not affect the criterion.

A nonzero Galois morphism $r\to r(1)$ would give a nonzero
morphism of Weil--Deligne representations. It would remain a
morphism after Frobenius semisimplification: an intertwiner of
Frobenius also intertwines its semisimple part, and the inertia
and monodromy conditions are unchanged. This contradicts the
genericity criterion.
\end{proof}

\begin{proposition}\label{prop:cohomology}
For the representation $r$ of a cuspidal newform of weight at
least $3$ and odd level dividing $N$,
\[
 \dim_E H^1(G_{\Q,S},\ad r)\leq3.
\]
\end{proposition}
\begin{proof}
Write $h^j$ for the $E$-dimension of a cohomology group.
Local Tate duality and the local Euler characteristic formulas
are used in their characteristic-zero form; see
\cite[Chapter~7]{NeukirchSchmidtWingberg08}.
Duality and Lemma~\ref{lem:local-invariants} give
$h^2(\Q_v,V)=0$ at every finite place. For $v\ne2$, the
local Euler characteristic formula yields
$h^1(\Q_v,V)=h^0(\Q_v,V)$. The unramified cohomology is
\[
 H^1_f(\Q_v,V)=
 V^{I_v}/(\Frob_v-1)V^{I_v}.
\]
The cokernel and kernel of an endomorphism of a finite-dimensional
vector space have the same dimension. Consequently this group
has dimension $h^0(\Q_v,V)$, and
\begin{equation}\label{eq:away-two}
 H^1(\Q_v,V)/H^1_f(\Q_v,V)=0\qquad(v\ne2).
\end{equation}

At $2$, the Euler characteristic formula instead gives
\[
 h^1(\Q_2,V)=h^0(\Q_2,V)+\dim_E V
            =h^0(\Q_2,V)+4.
\]
Write
$D_{\mathrm{dR}}(V)=(V\otimes_{\Q_2}B_{\mathrm{dR}})^{G_{\Q_2}}$
for the filtered de Rham module, where $B_{\mathrm{dR}}$ is
Fontaine's de Rham period field. The Bloch--Kato dimension formula
\cite[Corollary~3.8.4]{BlochKato90} gives
\begin{align*}
 h^1_f(\Q_2,V)
 &=h^0(\Q_2,V)+
   \dim_E\bigl(D_{\mathrm{dR}}(V)/\Fil^0D_{\mathrm{dR}}(V)\bigr)\\
 &=h^0(\Q_2,V)+1.
\end{align*}
Indeed the filtration degrees of the adjoint consist of two
zeros and two nonzero opposite integers, so exactly one degree
contributes to the indicated quotient. We have proved
\begin{equation}\label{eq:at-two}
 \dim_E H^1(\Q_2,V)/H^1_f(\Q_2,V)=3.
\end{equation}

Finally consider restriction followed by passage to these local
quotients:
\begin{equation}\label{eq:selmer-injection}
 H^1(G_{\Q,S},V)\longrightarrow
 \bigoplus_{v\in S}H^1(\Q_v,V)/H^1_f(\Q_v,V).
\end{equation}
Inflation identifies the source with a subspace of
$H^1(G_{\Q},V)$. Its classes are already unramified outside $S$,
so the kernel of \eqref{eq:selmer-injection} lies in
$H^1_f(\Q,V)$, which vanishes by
Lemma~\ref{lem:global-selmer}. Thus \eqref{eq:selmer-injection}
is injective. Equations~\eqref{eq:away-two} and
\eqref{eq:at-two} give a target of dimension $3$.
\end{proof}

\section{Dimension of the components}
\label{sec:conclusion}

\begin{proof}[Proof of Theorem~\ref{thm:main}]
Let $P$ be an arbitrary minimal prime of $A$. By
Corollary~\ref{cor:classical-point}, choose a classical cuspidal
point $\frakx=\ker\lambda$ of weight at least $3$ with
$P\subseteq\frakx$. Equation~\eqref{eq:hecke-tangent} and
Proposition~\ref{prop:cohomology} give
\[
 \dim A_{\frakx}\leq3.
\]
The quotient $A/P$ is a complete Noetherian local domain by
Proposition~\ref{prop:noetherian}, and is therefore catenary.
The classical quotient $A/\frakx$ is finite over $\Z_2$ and
contains $\Z_2$, so has dimension one. The dimension formula for
a catenary local domain yields
\[
 \dim A/P
 =\height(\frakx/P)+\dim A/\frakx
 \leq\dim A_{\frakx}+1
 \leq4.
\]
Theorem~\ref{thm:lower} gives the opposite inequality. Since $P$
was arbitrary, all irreducible components have dimension $4$.
\end{proof}

\bibliographystyle{amsplain}
\bibliography{references}
\end{document}
